% Part: incompleteness
% Chapter: incompleteness-provability
% Section: fixed-point-lemma

\documentclass[../../../include/open-logic-section]{subfiles}

\begin{document}

\olfileid{inc}{inp}{fix}

\olsection{స్థిరబిందు ఉపసిద్ధాంతం}

\begin{explain}
$\Th{T}$ అనేది $\Th{Q}$ను విస్తరిస్తే, ఏ \tetoken{సూత్రం}{!!{formula}}~$!B(x)$కైనా
$\Th{T} \Proves !A \liff !B(\gn{!A})$ అయ్యే \tetoken{ఒక వాక్యం}{!!a{sentence}}~$!A$ ఉంటుందని
స్థిరబిందు ఉపసిద్ధాంతం చెబుతుంది. అబద్ధిక వాక్యం విషయంలో, $!A$ అనేది
``$\gn{!A}$ అసత్యం''---అంటే $\Gn{!A}$ ఒక అసత్య \tetoken{వాక్యపు}{!!{sentence}} గ్యోడెల్
సంఖ్య---అనే ప్రకటనకు ($\Th{T}$లో నిరూపణీయంగా) సమానార్థకంగా ఉండాలని కోరుతాం.
నిరూపణ ఆలోచనను అర్థం చేసుకోవడానికి, $!A$ గురించి క్వైన్ ఇచ్చిన ఈ అనధికారిక
వివరణతో పోల్చడం ఉపయోగకరం: ``‘తన స్వంత ఉల్లేఖనాన్ని ముందు ఉంచితే అసత్యాన్ని
ఇస్తుంది’ అనే పదబంధం, తన స్వంత ఉల్లేఖనాన్ని ముందు ఉంచితే అసత్యాన్ని ఇస్తుంది.''
ఒక వ్యక్తీకరణను తీసుకుని, దానికి దాని స్వంత ఉల్లేఖనాన్ని ముందు ఉంచి వాక్యాన్ని
ఏర్పరచే క్రియను ఆ వ్యక్తీకరణను \emph{వికర్ణీకరించడం} అనవచ్చు; ఫలితాన్ని దాని
వికర్ణీకరణ అంటాం. కాబట్టి ‘తన స్వంత ఉల్లేఖనాన్ని ముందు ఉంచితే అసత్యాన్ని
ఇస్తుంది’ అనే పదబంధపు వికర్ణీకరణ ``‘తన స్వంత ఉల్లేఖనాన్ని ముందు ఉంచితే
అసత్యాన్ని ఇస్తుంది’ అనే పదబంధం, తన స్వంత ఉల్లేఖనాన్ని ముందు ఉంచితే
అసత్యాన్ని ఇస్తుంది.'' క్వైన్ అబద్ధిక వాక్యం ‘అసత్యాన్ని ఇస్తుంది’ అనే
పదబంధపు వికర్ణీకరణ కాదు; ‘తన స్వంత ఉల్లేఖనాన్ని ముందు ఉంచితే అసత్యాన్ని
ఇస్తుంది’ అనే పదబంధపు వికర్ణీకరణ అని గమనించండి. అంటే అబద్ధిక వాక్యాన్ని ఇచ్చే
వికర్ణీకృత లక్షణంలోనే వికర్ణీకరణ ఇమిడి ఉంది!{}

అంకగణిత భాషలో, ఒక స్వేచ్ఛా చరం గల \tetoken{ఒక సూత్రపు}{!!a{formula}} గ్యోడెల్ సంఖ్యను
గణించి, ఆ గ్యోడెల్ సంఖ్యకు చెందిన ప్రమాణ సంఖ్యాంకాన్ని స్వేచ్ఛా చరంలో
ప్రతిస్థాపించడం ద్వారా దాని ఉల్లేఖనాన్ని ఏర్పరుస్తాం. $!E(x)$ వికర్ణీకరణ
$!E(\num{n})$; ఇక్కడ $n = \Gn{!E(x)}$. (ఇక నుంచి $\num{\Gn{!E(x)}}$ను
$\gn{!E(x)}$గా సంక్షిప్తీకరిద్దాం.) కాబట్టి $!B(x)$ అంటే ``అసత్యం'' అయితే,
``తన స్వంత ఉల్లేఖనాన్ని ముందు ఉంచితే అసత్యాన్ని ఇస్తుంది'' అనేది
``తన వికర్ణీకరణ గ్యోడెల్ సంఖ్యకు అన్వయిస్తే అసత్యాన్ని ఇస్తుంది'' అవుతుంది.
\tetoken{సూత్రం}{!!{formula}} యొక్క గ్యోడెల్ సంఖ్య~$n$ ఇచ్చినప్పుడు దాని వికర్ణీకరణ గ్యోడెల్
సంఖ్యను గణించే $\fn{diag}(n)$ ప్రమేయానికి $\Obj{diag}$ అనే సంకేతం ఉంటే,
$!E(x)$ను $!B(\Obj{diag}(x))$గా రాయవచ్చు. అప్పుడు క్వైన్ అబద్ధిక వాక్యం
దాని వికర్ణీకరణ, అంటే $!E(\gn{!E(x)})$ లేదా
$!B(\Obj{diag}(\gn{!B(\Obj{diag}(x))}))$. ఇప్పుడు $!B(x)$ ఏ ఇతర లక్షణమైనా
కావచ్చు; ఇదే నిర్మాణం పనిచేస్తుంది. అసంపూర్ణతా సిద్ధాంతానికి $!B(x)$ను
``$x$ అనేది $\Th{T}$లో \tetoken{వ్యుత్పాదించదగినది}{!!{derivable}} కాదు''గా తీసుకుంటాం. అప్పుడు
$!E(x)$ అంటే ``తన వికర్ణీకరణ గ్యోడెల్ సంఖ్యకు అన్వయిస్తే $\Th{T}$లో
\tetoken{వ్యుత్పాదించదగినది}{!!{derivable}} కాని \tetoken{ఒక వాక్యాన్ని}{!!a{sentence}} ఇస్తుంది.''

దీనిని $\Th{T}$లో అధికారికీకరించడానికి $\fn{diag}$ను అధికారికీకరించే మార్గం
కనుగొనాలి. $\fn{diag}(n)$ గణనీయమైనది; నిజానికి ఆదిమ పునరావర్తనమైనది:
$n$ అనేది సూత్రం~$!E(x)$ గ్యోడెల్ సంఖ్య అయితే, $\fn{diag}(n)$ అనేది
$!E(\gn{!E(x)})$ గ్యోడెల్ సంఖ్యను ఇస్తుంది. ($\gn{!E(x)}$ అనేది
$!E(x)$ గ్యోడెల్ సంఖ్య యొక్క ప్రమాణ సంఖ్యాంకం, అంటే $\num{\Gn{!E(x)}}$ అని
గుర్తుంచుకోండి.) $\Obj{diag}$ అనేది $\fn{diag}$కు ప్రాతినిధ్యం వహించే
$\Th{T}$ ప్రమేయ సంకేతమైతే, $!A$ను
$!B(\Obj{diag}(\gn{!B(\Obj{diag}(x))}))$గా తీసుకోవచ్చు. గమనించండి:
\begin{align*}
\fn{diag}(\Gn{!B(\Obj{diag}(x))}) & =
\Gn{!B(\Obj{diag}(\gn{!B(\Obj{diag}(x))}))} \\
& = \Gn{!A}.
\end{align*}
$\Th{T}$ కింది దాన్ని \tetoken{నిగమించగలదని}{!!{derive}} అనుకుందాం:
\[
\Obj{diag}(\gn{!B(\Obj{diag}(x))}) = \gn{!A},
\]
అప్పుడు అది $!B(\Obj{diag}(\gn{!B(\Obj{diag}(x))}))
\liff !B(\gn{!A})$ను \tetoken{నిగమించగలదు}{!!{derive}}. కానీ నిర్వచనం ప్రకారం ఎడమవైపు
ఉన్నది~$!A$.

సాధారణంగా $\Obj{diag}$ అనేది $\Th{T}$లో ప్రమేయ సంకేతం కాదు; $\Th{Q}$లోని
సంకేతమైతే అసలే కాదు. కానీ $\fn{diag}$ గణనీయమైనది కాబట్టి, ఏదో సూత్రం
$!D_{\fn{diag}}(x,y)$ ద్వారా అది $\Th{Q}$లో \emph{ప్రాతినిధ్యం చేయదగినది}.
అందువల్ల $!B(\Obj{diag}(x))$కు బదులు
$\lexists[y][(!D_{\fn{diag}}(x,y) \land !B(y))]$ రాయవచ్చు. మిగతా నిరూపణ
పైన సూచించినట్లే సాగుతుంది; నిజానికి అది ఇప్పటికే $\Th{Q}$లోనే సాగుతుంది.
\end{explain}

\begin{lem}
\ollabel{lem:fixed-point} $x$ అనే ఒక స్వేచ్ఛా చరం గల ఏ సూత్రమైనా $!B(x)$గా
తీసుకోండి. అప్పుడు $\Th{Q} \Proves !A \liff !B(\gn{!A})$ అయ్యే వాక్యం
$!A$ ఉంటుంది.
\end{lem}


\begin{proof}
$!B(x)$ ఇచ్చినప్పుడు $!E(x)$ను
$\lexists[y][(!D_{\fn{diag}}(x,y) \land !B(y))]$గా, $!A$ను దాని
వికర్ణీకరణగా---అంటే $!E(\gn{!E(x)})$ అనే సూత్రంగా---తీసుకోండి.

$!D_{\fn{diag}}$ అనేది $\fn{diag}$కు ప్రాతినిధ్యం వహిస్తుంది, ఇంకా
$\fn{diag}(\Gn{!E(x)}) = \Gn{!A}$ కాబట్టి, $\Th{Q}$ కింది వాటిని
\tetoken{నిగమించగలదు}{!!{derive}}:
\begin{align}
  & !D_{\fn{diag}}(\gn{!E(x)}, \gn{!A}) \ollabel{repdiag1} \\
  & \lforall[y][(!D_{\fn{diag}}(\gn{!E(x)},y) \lif
  \eq[y][\gn{!A}])]. \ollabel{repdiag2}
\end{align}
ఇప్పుడు $\Th{Q} \Proves !A \liff !B(\gn{!A})$ అని చూపుదాం. తర్కాన్నీ,
$\Th{Q}$లో \tetoken{వ్యుత్పాదించదగిన}{!!{derivable}} వాస్తవాలనూ మాత్రమే వాడుతూ అనధికారికంగా
వాదిస్తాం.

మొదట $!A$, అంటే $!E(\gn{!E(x)})$ అని అనుకుందాం. $!E(x)$ నిర్వచనానికి
తిరిగి వెళితే, $!E(\gn{!E(x)})$ అనేది ఇదే:
\[
\lexists[y][(!D_{\fn{diag}}(\gn{!E(x)},y) \land !B(y))].
\]
అలాంటి~$y$ను పరిశీలించండి. $!D_{\fn{diag}}(\gn{!E(x)},y)$ కాబట్టి,
\olref{repdiag2} ప్రకారం $y = \gn{!A}$. అందువల్ల $!B(y)$ నుంచి
$!B(\gn{!A})$ వస్తుంది.

ఇప్పుడు $!B(\gn{!A})$ అని అనుకుందాం. \olref{repdiag1} ప్రకారం
\begin{align*}
& !D_{\fn{diag}}(\gn{!E(x)}, \gn{!A}) \land !B(\gn{!A}).
\intertext{దీనినుంచి}
& \lexists[y][(!D_{\fn{diag}}(\gn{!E(x)},y) \land !B(y))].
\end{align*}
కానీ ఇది $!E(\gn{!E(x)})$, అంటే~$!A$.
\end{proof}

\begin{digress}
దీనిని గణనీయతా సిద్ధాంతంలోని స్థిరబిందు ఉపసిద్ధాంతపు నిరూపణతో పోల్చాలి.
ఇక్కడ ఒక \emph{ప్రకటనను} దాని పరంగానే నిర్వచించాలి; అక్కడ ఒక
\emph{ప్రమేయాన్ని} దాని పరంగానే నిర్వచించాల్సి వచ్చింది. ఈ తేడా తప్పిస్తే,
రెండింటిలోనూ నిజంగా ఒకే ఆలోచన ఉంది.
\end{digress}

\begin{prob}
అన్ని \tetoken{వాక్యాలు}{!!{sentence}s}~$!B$కు $\Th{Q} \Proves !B \liff !A(\gn{!B})$
అయితే, \tetoken{ఒక సూత్రాన్ని}{!!^a{formula}}~$!A(x)$ \emph{సత్య నిర్వచనం} అంటాం. స్థిరబిందు
ఉపసిద్ధాంతాన్ని వాడి ఏ \tetoken{సూత్రమూ}{!!{formula}} సత్య నిర్వచనం కాదని చూపండి.
\end{prob}

\end{document}
